\section{Conditioning, safe event bounds, and continuous transitions}
\label{sec:conditioning}

\subsection{What an upper probability assertion means without an event name}

For a general $L^0$ variable $Y$, the expression $\{Y\ge a\}$ need not
come with a regular event name. A numerical bound on that event should not
be mistaken for a procedure computing its probability. There is nevertheless
an entirely defined observation that expresses the useful estimate.
For rational $0<\eta<a$, put
\[
 r_{a,\eta}(y)
 =\min\left(1,\max\left(0,\frac{y-a+\eta}{\eta}\right)\right).
\]
This bounded Lipschitz function is zero below $a-\eta$ and one at and above
$a$. Its expectation always exists for an $L^0$ name. For $Y\ge0$ in $L^1$,
\[
 \mathbb E r_{a,\eta}(Y)\le\frac{\mathbb EY}{a-\eta}.
\]
The pointwise finite inequality follows by splitting at $a-\eta$ and $a$;
in the transition interval the right side is at least one. For $X\in L^2$,
apply the analogous bound to $r_{a,\eta}(|X|)$ to get
$\mathbb E r_{a,\eta}(|X|)\le\mathbb E X^2/(a-\eta)^2$.

In this report an \emph{upper-envelope conclusion} refers to this family of
defined inequalities with explicit rational slack. For any positive target
slack, a small enough rational $\eta$ is selected from $a$ and the moment
bound by bounded arithmetic. No unspecified outer-measure real is being
postulated. If the hard event has a boundary certificate, the same
observations converge to its indicator with a PR rate and give the usual
probability bound. In standard measure semantics they majorize the hard
event whether or not that event has a computable probability. The same
construction applies to finite maxima, which are continuous functions of
finitely many coordinates.

\subsection{Bayes updating with a certified evidence bound}

Let $\ell\ge0$ be a represented likelihood bounded by $L$, and suppose
$Z=\mathbb E\ell\ge z_0>0$ is certified. On represented bounded
observables define
\[
 \mathbb E_\ell h=\frac{\mathbb E(h\ell)}{Z},\qquad
 \mathbb P_\ell(A)=\frac{\mathbb E(\mathbf1_A\ell)}{Z}.
\]
Bounded multiplication and the reciprocal construction make these regular
real names. Positivity, normalization, and finite additivity follow from the
corresponding expectation identities. Moreover
\[
 |\mathbb P_\ell(A)-\mathbb P_\ell(B)|
 \le (L/z_0)\,\mathbb P(A\triangle B),
\]
so the construction respects completed event equality and transfers
countable-union moduli with the indicated factor.

The numerical stability estimate is explicit. Suppose $|h|\le M$, the
numerator $a=\mathbb E(h\ell)$ is approximated by $a'$ to error
$\delta_a$, and $Z$ by $Z'$ to error $\delta_Z\le z_0/2$. Since
$|a|\le MZ$ and $Z'\ge z_0/2$,
\begin{equation}\label{eq:bayes-budget}
 \left|\frac{a'}{Z'}-\frac aZ\right|
 \le\frac{2(\delta_a+M\delta_Z)}{z_0}.
\end{equation}
This follows by adding and subtracting $a/Z'$ and bounding the two terms.
Conditioning on an event is the case $\ell=\mathbf1_B$. Bayes' identity is
then an equality of intersection probabilities, with every denominator's
positive margin retained. No search for a positive lower bound is inferred
from a bare negation of zero.

For a joint density $f(x,y)$, a pointwise conditional density can be
constructed on a compact observation region if its section integrals have
uniform quadrature and tail data and the marginal
$m(y)=\int f(x,y)\,dx$ has a certified lower bound there. Then
$f(x,y)/m(y)$ is a normalized density in $x$. Finite Fubini, with the
supplied uniform tails, proves the conditional integration identity
\[
 \int_B\left(\int_A\frac{f(x,y)}{m(y)}\,dx\right)m(y)\,dy
   =\int_{A\times B}f(x,y)\,dx\,dy
\]
for the represented integrable tests or regular events at issue. The
Gaussian example below supplies all these data explicitly. General
conditional distributions are not automatically computable from a computable
joint law; counterexamples in the computability literature confirm that
these restrictions are substantive \cite{ackerman}.

\subsection{A finite stopping theorem and a Borel--Cantelli fragment}

Let $M_j$ be an integrable martingale for the first-$j$-bit filtration, with
the conditional identities presented as above. Let $\tau\le N$ be a
stopping rule whose decision $\{\tau>j\}$ is a certified event depending
only on the first $j$ observations. In the simplest internal implementation
these are finite cylinder decisions. On a common finite table,
\[
 M_\tau=M_0+\sum_{j<N}\mathbf1_{\{\tau>j\}}(M_{j+1}-M_j).
\]
Each summand has expectation zero by the conditional identity and past
measurability. Hence $\mathbb E M_\tau=\mathbb E M_0$. Bounded-multiplier
closure justifies the represented products. This is a bounded stopping
theorem; an unbounded stopping time requires additional uniform-integrability
and approximation data before a limit can be taken.

For a jointly PR event family $A_j$, suppose the finite tail sums of its
probabilities are at most $2^{-k}$ after a supplied index $C(k)$.
The admitted tail union $B_k=\bigcup_{j\ge C(k)}A_j$ has mass at most
$2^{-k}$, by finite union bounds and the explicit residual modulus. Every
later $A_j$ is included in $B_k$ in the event-algebra order. These small
tail events are an effective Borel--Cantelli certificate. Its familiar
interpretation is that only finitely many $A_j$ occur almost surely. The
arithmetic assertion is the positive inclusion and small-mass certificate,
not a nonconstructive inference from a logical negation of convergence.

\subsection{Continuous-state Markov chains with supplied transition data}

A useful transition presentation supplies a function $\Psi(x,u)$ taking a
state and an independent uniform input, a finite approximation scheme, and
stability and tail bounds. For example, a coupling estimate
\[
 \mathbb E\min(1,|\Psi(x,U)-\Psi(y,U)|)
 \le C\min(1,|x-y|)
\]
with rational $C\ge1$ propagates $L^0$ approximation errors.
Using independent bit rows, set $X_{j+1}=\Psi(X_j,U_j)$. This constructs
every finite trajectory jointly. If each numerical transition has error
at most $\delta_j$, the finite recurrence
\[
 e_{j+1}\le Ce_j+\delta_j
 \quad\Longrightarrow\quad
 e_n\le C^ne_0+\sum_{j<n}C^{n-1-j}\delta_j
\]
is a quantifier-free induction. Choosing all input errors at most
$2^{-k}/[(n+1)C^n]$ suffices at horizon $n$. This is a concrete PR
construction of a continuous-state process at finite horizons. An arbitrary
family of finite-dimensional laws, without coherent realizers and effective
tightness, is not a substitute for these data.

For an explicit example take $X_0=x_0$ and
\[
 X_{j+1}=aX_j+b+\sigma G_j,
 \qquad |a|\le r<1,\quad\sigma\ge\sigma_0>0,
\]
where $G_j$ are independent standard normals and all parameters have
represented bounds. Expansion and Gaussian convolution show, for $n\ge1$,
that $X_n$ is normal with
\[
 m_n=a^nx_0+b\sum_{j<n}a^j,\qquad
 v_n=\sigma^2\sum_{j<n}a^{2j}.
\]
These approach $m=b/(1-a)$ and $v=\sigma^2/(1-a^2)$ with bounds
\[
 |m_n-m|\le r^n|x_0-m|,\qquad
 |v_n-v|\le\frac{\sigma^2r^{2n}}{1-r^2}.
\]
The identities $am+b=m$ and $a^2v+\sigma^2=v$, together with the
Gaussian transition identity, show that $N(m,v)$ is invariant under one
transition. This identity is valid also when $a=0$ and does not require
deciding whether the represented parameter $a$ vanishes.
Couple the two normal laws as $m_n+\sqrt{v_n}G$ and $m+\sqrt vG$.
Since $v_n,v\ge\sigma_0^2$ and $\mathbb E|G|\le1$,
\[
 \mathbb E|X_n-X_\infty|
 \le r^n|x_0-m|
       +\frac{\sigma^2r^{2n}}{2\sigma_0(1-r^2)}.
\]
Here the displayed coupling compares laws and is not a claim that the
original noisy trajectory converges to a single random variable.
If $r\le1-q$ for rational $q>0$, the elementary bound
$(1-q)^n\le1/(1+nq)$ yields a PR precision schedule. Thus stationary-law
convergence is available without an abstract ergodic theorem. Continuous-time
path constructions and stochastic integration require further modules and
are not claimed in this package.
